\documentclass[11pt, a4paper]{article}
\usepackage[utf8]{inputenc}
\usepackage[T1]{fontenc}
\usepackage[spanish]{babel}
\usepackage[a4paper, margin=2.5cm]{geometry}
\usepackage{amsmath, amssymb}
\usepackage{booktabs}
\usepackage{enumitem}
\setlist[itemize]{label=-}
\usepackage{hyperref}

\title{\textbf{Clasificación Modular y Análisis de Admisibilidad de Primos Gemelos mediante un Molde Séxtuple}}
\author{Cuaderno de Investigación: Primos Gemelos}
\date{}

\begin{document}

\maketitle

\begin{abstract}
En este documento se formalizan los resultados obtenidos sobre la distribución de los números primos gemelos mayores que $5$ mediante aritmética modular y conjuntos admisibles. Se demuestra que todo par de primos gemelos $(p, p+2)$ con $p > 5$ pertenece necesariamente a una de tres clases residuales módulo $30$, las cuales pueden agruparse dentro de un único molde admisible de seis elementos generado por la progresión aritmética $x = 11 + 30n$. Se incluye además la verificación computacional explícita para el rango $0 \le n \le 9$ y un análisis de simetría frente a los divisores primos $p \ge 7$.
\end{abstract}

\section{Concepto de Admisibilidad}

Sea $H = \{h_1, h_2, \dots, h_k\}$ un conjunto finito de enteros distintos. Se dice que $H$ es \textbf{admisible} si, para todo número primo $p$, existe al menos un entero $x$ tal que ninguno de los elementos del conjunto desplazado
\[
x + H = \{x + h_1, x + h_2, \dots, x + h_k\}
\]
es divisible por $p$. Equivalentemente, los elementos de $H$ no cubren todas las clases de congruencia módulo $p$ para ningún primo $p$.

\begin{itemize}
    \item \textbf{Ejemplo admisible:} El conjunto $\{0, 4, 6, 10\}$ es admisible. Para $p = 2$ y $p = 3$, basta elegir $x = 1$, obteniendo $\{1, 5, 7, 11\}$, ninguno de cuyos elementos es divisible por $2$ ni por $3$. Para $p \ge 5$, al haber solo cuatro elementos, siempre queda al menos un residuo libre.
    \item \textbf{Ejemplo no admisible:} El conjunto $\{0, 2, 4, 8\}$ no es admisible porque, independientemente del valor de $x$, el subconjunto $\{x, x+2, x+4\}$ cubre los tres residuos posibles módulo $3$, garantizando siempre la aparición de un múltiplo de $3$.
\end{itemize}

\section{Partición Módulo 30 y el Molde Séxtuple}

Todo número primo mayor que $5$ es coprimo con el primorial $W = 2 \times 3 \times 5 = 30$. En consecuencia, cualquier primo $p > 5$ debe dejar uno de los ocho residuos coprimos con $30$:
\[
p \equiv 1, 7, 11, 13, 17, 19, 23, 29 \pmod{30}.
\]

Si exigimos que $(p, p+2)$ sea un par de primos gemelos con $p > 5$, únicamente existen tres parejas de residuos consecutivos separados por una distancia de $2$:
\begin{enumerate}
    \item \textbf{Pareja A:} $(30n + 11,\; 30n + 13)$
    \item \textbf{Pareja B:} $(30n + 17,\; 30n + 19)$
    \item \textbf{Pareja C:} $(30n + 29,\; 30n + 31)$
\end{enumerate}

Estas tres familias pueden integrarse en un único molde admisible de seis posiciones. Definimos el desplazamiento base $H = \{0, 2, 6, 8, 18, 20\}$ y el punto de partida $x = 11 + 30n$ para $n \in \{0, 1, 2, \dots\}$. De este modo, cada valor de $n$ genera la séxtupla:
\[
S(n) = \{x,\; x+2,\; x+6,\; x+8,\; x+18,\; x+20\}.
\]

Por construcción, ningún elemento de $S(n)$ es divisible por $2$, $3$ ni $5$, y todo par de primos gemelos mayor que $(5, 7)$ aparece necesariamente en una de las tres parejas internas de $S(n)$ para algún $n \ge 0$.

\section{Verificación Explícita para $0 \le n \le 9$}

La Tabla \ref{tab:sextuplas} muestra las diez primeras séxtuplas generadas por $x = 11 + 30n$, indicando los números primos en negrita y la cantidad total de primos por fila.

\begin{table}[h!]
\centering
\caption{Séxtuplas generadas por $S(n)$ en el intervalo $0 \le n \le 9$.}
\label{tab:sextuplas}
\vspace{0.2cm}
\begin{tabular}{@{}c c l c l@{}}
\toprule
$n$ & $x$ & Séxtupla $S(n)$ & Primos & Factorización de compuestos \\
\midrule
0 & 11  & $\{\mathbf{11}, \mathbf{13}, \mathbf{17}, \mathbf{19}, \mathbf{29}, \mathbf{31}\}$ & 6 & Ninguno \\
1 & 41  & $\{\mathbf{41}, \mathbf{43}, \mathbf{47}, 49, \mathbf{59}, \mathbf{61}\}$ & 5 & $49 = 7^2$ \\
2 & 71  & $\{\mathbf{71}, \mathbf{73}, 77, \mathbf{79}, \mathbf{89}, 91\}$ & 4 & $77 = 7 \cdot 11,\; 91 = 7 \cdot 13$ \\
3 & 101 & $\{\mathbf{101}, \mathbf{103}, \mathbf{107}, \mathbf{109}, 119, 121\}$ & 4 & $119 = 7 \cdot 17,\; 121 = 11^2$ \\
4 & 131 & $\{\mathbf{131}, 133, \mathbf{137}, \mathbf{139}, \mathbf{149}, \mathbf{151}\}$ & 5 & $133 = 7 \cdot 19$ \\
5 & 161 & $\{161, \mathbf{163}, \mathbf{167}, 169, \mathbf{179}, \mathbf{181}\}$ & 4 & $161 = 7 \cdot 23,\; 169 = 13^2$ \\
6 & 191 & $\{\mathbf{191}, \mathbf{193}, \mathbf{197}, \mathbf{199}, 209, \mathbf{211}\}$ & 5 & $209 = 11 \cdot 19$ \\
7 & 221 & $\{221, \mathbf{223}, \mathbf{227}, \mathbf{229}, \mathbf{239}, \mathbf{241}\}$ & 5 & $221 = 13 \cdot 17$ \\
8 & 251 & $\{\mathbf{251}, 253, \mathbf{257}, 259, \mathbf{269}, \mathbf{271}\}$ & 4 & $253 = 11 \cdot 23,\; 259 = 7 \cdot 37$ \\
9 & 281 & $\{\mathbf{281}, \mathbf{283}, 287, 289, 299, 301\}$ & 2 & $287 = 7 \cdot 41,\; 289 = 17^2,\; 299 = 13 \cdot 23,\; 301 = 7 \cdot 43$ \\
\bottomrule
\end{tabular}
\end{table}

\section{Simetría Frente a Primos $p \ge 7$ y Séxtuplas Completas}

Al analizar la distribución de los pares gemelos dentro de $S(n)$ para $0 \le n \le 9$ (todos los primos menores que $301$ excepto $(3,5)$ y $(5,7)$), se observa un reparto equilibrado entre las tres familias:
\begin{itemize}
    \item \textbf{Pareja A:} $6$ pares gemelos: $(11,13)$, $(41,43)$, $(71,73)$, $(101,103)$, $(191,193)$, $(281,283)$.
    \item \textbf{Pareja B:} $5$ pares gemelos: $(17,19)$, $(107,109)$, $(137,139)$, $(197,199)$, $(227,229)$.
    \item \textbf{Pareja C:} $6$ pares gemelos: $(29,31)$, $(59,61)$, $(149,151)$, $(179,181)$, $(239,241)$, $(269,271)$.
\end{itemize}

Respecto a la aparición de séxtuplas donde los seis elementos sean primos simultáneamente, es condición necesaria que ningún elemento sea divisible por $7$. Dado que $30 \equiv 2 \pmod 7$, la condición $7 \nmid S(n)$ exige que
\[
n \equiv 0 \pmod 7.
\]
En el intervalo $1 \le n \le 1000$, ningún múltiplo de $7$ logra producir seis primos simultáneos debido a obstrucciones de primos mayores ($11, 13, 17, 19, 23$). El siguiente índice después de $n = 0$ donde los seis elementos vuelven a ser primos ocurre en $n = 553$, correspondiente al punto de partida $x = 16601$:
\[
S(553) = \{16601,\; 16603,\; 16607,\; 16609,\; 16619,\; 16621\}.
\]

\end{document}